\documentclass[a4paper,10pt]{article}

\usepackage[utf8]{inputenc}
\usepackage[T1]{fontenc}
\usepackage[german]{babel}
\usepackage{geometry}
\usepackage{xcolor}
\usepackage{titlesec}
\usepackage{hyperref}
\usepackage{fontawesome5}
\usepackage{charter}

% Configuration
\geometry{left=2.5cm, top=2cm, right=2.5cm, bottom=2cm}
\definecolor{primarycolor}{RGB}{0, 51, 102}
\hypersetup{colorlinks=true, urlcolor=primarycolor}

% Custom commands
\newcommand{\contactitem}[2]{\small#1 \ #2 \quad}

\begin{document}

% Header
\begin{flushleft}
    {\Huge \textbf{\color{primarycolor} Kartavya Niraj Dikshit}} \\
    \vspace{4pt}
    \small
    Halbmondstraße 7, 91054 Erlangen \\
    \faPhone* (+49) 15510940829 \\
    \faEnvelope \href{mailto:kartavya.dikshit@gmail.com}{kartavya.dikshit@gmail.com} \\
    \faLinkedin \href{http://www.linkedin.com/in/kartavya-dikshit}{linkedin.com/in/kartavya-dikshit} \\
    \faGithub \href{https://github.com/KartavyaDikshit}{github.com/KartavyaDikshit}
\end{flushleft}

\vspace{20pt}

\begin{flushleft}
    Siemens AG \\
    Smart Infrastructure \\
    Recruiting Team \\
    Werner-von-Siemens-Straße 48 \\
    91052 Erlangen
\end{flushleft}

\vspace{20pt}

\begin{flushright}
    Erlangen, \today
\end{flushright}

\vspace{10pt}

\noindent \textbf{\large Bewerbung als Werkstudent (w/m/d) AI-gestütztes Kosten- und Value Engineering – Softwareentwicklung (Job ID: 496184)}

\vspace{20pt}

Sehr geehrte Damen und Herren,

als Masterstudent der Data Science mit dem Schwerpunkt Künstliche Intelligenz an der Friedrich-Alexander-Universität Erlangen-Nürnberg (FAU) verfolge ich mit großem Interesse, wie Siemens Smart Infrastructure die Grenzen des technisch Machbaren durch Digitalisierung verschiebt. Die Verknüpfung von Softwareentwicklung mit datengesteuerter wirtschaftlicher Analyse im Kosten- und Value Engineering ist genau die Schnittstelle, an der ich meine akademischen Kenntnisse und praktischen Erfahrungen wertbringend einsetzen möchte.

Mein aktuelles Studium an der FAU vertieft mein Verständnis für fortgeschrittene Machine-Learning-Methoden und statistische Modellierung. Diese theoretische Basis habe ich bereits in anspruchsvollen Projekten angewendet. So entwickelte ich eine \textit{Predictive Industrial Cost-Variance Engine}, die mithilfe von Gradient Boosted Decision Trees (GBDT) historische Kostenstrukturen analysiert und Prognosefehler in der Beschaffung signifikant reduzierte. Diese Erfahrung in der Aufbereitung komplexer Datensätze für wirtschaftliche Analysen korrespondiert direkt mit den Anforderungen Ihrer Ausschreibung zur Optimierung von KI-Algorithmen im Controlling.

Darüber hinaus bringe ich fundierte Kenntnisse in der Softwareentwicklung mit Python, Java und C# mit. In meiner bisherigen Tätigkeit als AI/ML-Engineer habe ich skalierbare RAG-Systeme (Retrieval-Augmented Generation) implementiert, um technische Dokumentationen automatisiert zu klassifizieren – ein Ansatz, den ich für die Wertanalyse (Value Engineering) bei Siemens adaptieren möchte, um Effizienzpotenziale in Engineering-Spezifikationen schneller zu identifizieren. Die Zusammenarbeit mit interdisziplinären Teams ist für mich durch meine Erfahrungen in agilen R&D-Umgebungen, unter anderem bei Icertis, gelebte Praxis.

Ich bin davon überzeugt, dass ich durch meine Kombination aus technischer Tiefe in der KI-Entwicklung und einem ausgeprägten Verständnis für wirtschaftliche Zusammenhänge einen entscheidenden Beitrag zur Zukunftsfähigkeit Ihrer Engineering-Lösungen leisten kann. Mein Wohnsitz in Erlangen ermöglicht mir zudem eine flexible und langfristige Zusammenarbeit von 15 bis 20 Stunden pro Woche.

Gerne möchte ich Sie in einem persönlichen Gespräch von meiner Motivation und meinen Fähigkeiten überzeugen.

Mit freundlichen Grüßen

\vspace{20pt}

Kartavya Niraj Dikshit

\end{document}